% Part: incompleteness
% Chapter: representability-in-q
% Section: prim-rec

\documentclass[../../../include/open-logic-section]{subfiles}

\begin{document}

\olfileid{inc}{req}{pri}
\olsection{محاكاة العودية البدائية}

يمكننا الآن أن نثبت أن التعريف بالعودية البدائية يمكن «محاكاته»
بالتصغير المنتظم باستعمال دالة بيتا. افترض أن لدينا $f(\vec x)$ و
$g(\vec x,y,z)$. عندئذ تكون الدالة~$h(\vec x,y)$ المعرّفة من $f$
و$g$ بالعودية البدائية هي
\begin{align*}
h(\vec x, 0) & =  f(\vec x) \\
h(\vec x, y+1) & =  g(\vec x, y, h(\vec x, y)).
\end{align*}
وعلينا أن نثبت إمكان تعريف $h$ من $f$ و$g$ باستعمال التركيب والتصغير
المنتظم وحدهما، مع الدوال الأساسية والدوال المعرّفة منها بالتركيب
والتصغير المنتظم (مثل~$\beta$).

\begin{lem}
\ollabel{lem:prim-rec}
إذا أمكن تعريف $h$ من $f$ و$g$ باستعمال العودية البدائية، أمكن
تعريفها من $f$ و$g$ والدوال $\Zero$ و$\Succ$ و$\Proj{n}{i}$ و$\Add$
و$\Mult$ و$\Char{=}$، باستعمال التركيب والتصغير المنتظم.
\end{lem}

\begin{proof}
عرّف أولًا دالة مساعدة $\hat h(\vec x,y)$ تعيد أصغر عدد~$d$ بحيث يرمّز
$d$ متتالية تحقق
\begin{enumerate}
\item $(d)_0=f(\vec x)$، و
\item لكل $i<y$، يكون $(d)_{i+1}=g(\vec x,i,(d)_i)$،
\end{enumerate}
حيث يختصر $(d)_i$ الآن $\beta(d,i)$. وبعبارة أخرى، تعيد $\hat h$
المتتالية
$\tuple{h(\vec x,0),h(\vec x,1),\dots,h(\vec x,y)}$. ويمكننا كتابة
$\hat h$ بالصيغة
\[
\hat h(\vec x, y) = \umin{d}{(\beta(d,0) = f(\vec x) \land \bforall{i <
  y}{\beta(d,i+1) = g(\vec x, i,\beta(d,i)})}.
\]
لاحظ أنه لا حاجة هنا إلى العودية البدائية، بل إلى التصغير وحده.
والدالة التي نصغّرها منتظمة بفضل لمة دالة بيتا
\olref[bet]{lem:beta}.

لكن لدينا الآن
\[
h(\vec x, y) = \beta(\hat h(\vec x, y), y),
\]
ومن ثم يمكن تعريف $h$ من الدوال الأساسية باستعمال التركيب والتصغير
المنتظم وحدهما.
\end{proof}

\end{document}
